\section{王鹤的学术路径：从视频到合成数据}

上一章讲学术群体，本章讲王鹤个人路径，因为他的博士项目几乎预演了后来具身智能的两条主线：从人类视频学习交互关系，以及用合成数据解决真实标注稀缺。第一项工作是从人类视频中学习多步人与物体交互，生成动画或驱动一个简单的智能杯子 demo。它本质上要学习状态、动作、物体状态变化和下一步可能动作之间的因果关系。

这件事在 2016 年很早。那时没有今天的多模态大模型，也没有成熟的 VLA 数据。要知道杯子空/满、瓶盖开/关、人手空/握物、动作片段边界，都需要拆成多个模型、手工标注和复杂系统。王鹤后来总结，虽然完全从视频学习世界模型在今天仍未成为最能推进具身智能的技术，但这个项目让他很早意识到：真正的问题不是识别，而是从交互中学出世界状态变化。

\begin{figure}[H]
\centering
\includegraphics[width=0.96\textwidth]{academic-path.png}
\caption{王鹤学术路径：从人类视频、世界模型到家庭机器人研究目标。自制概念图，依据 00:41:15--01:25:08 对谈内容整理。}
\end{figure}

\begin{knowledgebox}{读图：两个博士项目对应两条具身主线}
第一条是从视频中理解人和物体的多步因果交互，接近今天的 world model 问题；第二条是类别级物体位姿估计，用合成/混合现实数据解决真实标注不足。读图时要把它们看成同一目标的两个侧面：理解交互状态，以及低成本获得可训练数据。
\end{knowledgebox}

\subsection{从物理建模到 AI 研究}

本节解释为什么王鹤能把复杂交互抽象成模型。他本科和早期博士背景偏物理、器件和数学建模，习惯从实验数据中拟合模型，再用模型预测新 case。后来他发现半导体 clean room 的实验迭代慢、控制差，与自己的思维节奏不匹配，于是转向 AI、图形学和物理交互。这段经历说明，具身智能不只是写代码，还需要理解物体状态、动作因果和物理变化。

博士第一项项目最难的不是某个网络结构，而是 formulate research problem：哪些状态需要被抽取，哪些动作改变哪些物体状态，如何把这些变化串成时序因果模型。王鹤提到自己在还不知道 Perception-Action Loop 这个流行表达时，就给导师画过 state-action-change world-state 的图。这正是后来的具身智能底层循环。

\begin{knowledgebox}{术语消化：交互建模}
\begin{tabularx}{\textwidth}{p{0.20\textwidth}p{0.34\textwidth}X}
\toprule
术语 & 含义 & 在本期中的作用 \\
\midrule
Object State & 物体状态，如空/满、开/关、位置/朝向 & 决定下一步动作是否可行。 \\
Action Grammar & 用规则描述动作顺序的传统方法 & 可解释但泛化和数据驱动能力有限。 \\
World Model & 预测动作如何改变世界状态的模型 & 具身智能长期目标之一。 \\
Perception-Action Loop & 感知、行动、反馈的循环 & 把静态识别变成闭环智能。 \\
\bottomrule
\end{tabularx}
\end{knowledgebox}

\subsection{类别级位姿估计与合成数据}

第二项工作更接近当下具身数据路线。传统 6D pose estimation 往往针对一个已知实例：先给物体定义坐标原点和标准方向，再预测它相对坐标系的旋转和平移。王鹤的问题是类别级位姿估计：面对从未见过的马克杯，模型能否知道杯口朝上、杯柄朝右，而不需要每个杯子都扫描建模、逐个定义坐标系。

真实数据几乎不存在。互联网图片没有六维位姿标注，博士生也不可能拍摄并标注无限多马克杯。于是他采用图形学和混合现实数据：去宜家拍真实桌面背景和深度，把数字马克杯渲染到真实背景上，自动获得位姿标签，生成几十万张训练图。这个方法和今天的合成数据、Sim2Real 思想高度一致：真实背景提供接地，虚拟前景提供可控标签，最终在真实世界评估。

\begin{importantbox}{合成数据的早期版本}
类别级位姿估计说明，合成数据不是今天才出现的潮词。只要真实数据标注不可规模化，研究者就会用图形学、仿真或混合现实来制造可控训练样本。今天的机器人合成数据只是把这个思想扩展到更多物体、动作、场景和物理交互。
\end{importantbox}

\subsection{为什么坚持家庭机器人}

本节把个人路径转到战略选择。2020 年前后，李开复曾建议王鹤把三维视觉和位姿估计能力用于自动驾驶、LiDAR 和车的位姿估计；王鹤则坚持做 home robot，因为他认为家庭场景的交互空间最大，也最接近“物理世界通用智能体”。这在当时很不合时宜：家用机器人被认为还有很远，国内也几乎没有盟军。

回国后，王鹤在北大和智源推动具身智能方向，写文章介绍 embodied AI，和卢策武等少数研究者一起把这个概念带入国内学术与产业讨论。这里的教学重点不是个人传奇，而是早期方向选择的结构：如果一个研究者相信 next wave，就必须忍受一段没有共识、没有盟军、没有资本热度的时期。

\subsection{本章小结}

王鹤的学术路径把具身智能的两个关键问题提前暴露出来：第一，智能要理解动作如何改变世界；第二，真实数据缺口需要合成/混合现实数据补足。家庭机器人目标则把这些研究问题推向更大的物理世界应用。

